\section{От локального разрешённого сдвига к макроскопической ползучести}

Результаты раздела~3 ограничивают то, что деформированное включение
способно сделать локально. Чтобы связать их с макроскопическим наблюдением,
вопрос обращается: какое напряжение вокруг включений потребовалось бы, чтобы
повысить скорость ползучести на измеренные 25\%? Оценка, отвечающая на этот
вопрос, складывается из трёх составляющих.

Первая --- форма поля напряжений. Вне компактной частицы радиуса $a$,
сцепленной с матрицей и деформированной, упругое решение Эшелби [7] даёт
напряжение, спадающее с расстоянием $r$ от центра частицы как $(a/r)^{3}$.
Разрешённое касательное напряжение --- та часть напряжения, которая действует
толкающая дислокацию (раздел~2.4), и потому единственная
движет дислокацию, --- записывается в этой скалярной форме как
$\tau(r) = \tau_m (a/r)^{3}$, где $\tau_m$ --- его значение на поверхности
частицы. В разбавленной дисперсии с объёмной долей включений $f$ доля объёма
матрицы, в которой $|\tau|$ превышает заданный уровень $s$, равна
$F(>s) = f\,(\tau_m/s - 1)$ при $0 < s < \tau_m$ (объём самой частицы
вычтен), а распределение уровней напряжения по объёму матрицы задаётся
плотностью $p(s) = -\mathrm{d}F/\mathrm{d}s = f\,\tau_m/s^{2}$.

Вторая --- отклик скольжения дислокаций на напряжение. При комнатной
температуре дислокации преодолевают препятствия термически активированными
скачками, и локальное напряжение $\tau$ изменяет частоту скачков в
$\exp(\pm V^{*}\tau/kT)$ раз. Здесь $V^{*}$ --- активационный объём раздела~3.3 [18]. Поскольку напряжение вокруг
частицы принимает оба знака, оба направления учитываются с равным весом, и
скорость оказывается пропорциональной $\cosh(V^{*}\tau/kT)$.

Третья --- усреднение. Скорость ползучести сплава есть скорость, усреднённая
по объёму матрицы, поэтому $\cosh(V^{*}\tau/kT) - 1$ интегрируется с
плотностью $p(s)$. Относительное усиление скорости ползучести составляет
\begin{equation}
\frac{\langle\dot\varepsilon\rangle}{\dot\varepsilon_0} - 1
 = f\,x_0\!\int_{0}^{x_0}\!\frac{\cosh u - 1}{u^{2}}\,\mathrm{d}u
 = f\left[x_0\,\mathrm{Shi}(x_0) - (\cosh x_0 - 1)\right],
\qquad x_0 = \frac{V^{*}\tau_m}{kT},
\end{equation}
где $\mathrm{Shi}(x) = \int_0^x \sinh u\,\mathrm{d}u/u$ --- гиперболический
интегральный синус. Интеграл берётся до $s = 0$, то есть предполагается, что
напряжённые оболочки соседних частиц не перекрываются; обрезание его там, где
они перекрылись бы (при $u_{\min} = f x_0$, когда оболочка одной частицы
достигает соседних), понижает предсказанное усиление с 0{,}2500 до 0{,}2499 и
повышает требуемое $\tau_m$ примерно на $0{,}01\%$ при каждом $V^{*}$
табл.~\ref{tab:bridge}, так что приведённые значения сохраняются.
Соотношение~(2) --- модельная оценка чувствительности, а не извлечение
межфазного напряжения из эксперимента: оно описывает поле одной амплитудой
$\tau_m$, учитывает положительные и отрицательные напряжения с равным весом и
берёт пространственную статистику поля $(a/r)^{3}$. Существенно, что
$\tau_m$ --- амплитуда разрешённого сдвига на поверхности частицы, а не
полное межфазное напряжение, из которого в любой плоскости скольжения
действует лишь часть.

В оценку входят следующие величины. Температура --- 300~К. Содержание включений в сплаве
0{,}35~масс.\%, что отвечает объёмной доле $f = 0{,}00246$, в пределах
0{,}1--0{,}3\%, указанных в [5]; табл.~\ref{tab:bridge} рассчитана при этом
значении.
Активационный объём взят в интервале $V^{*} = 19$--$142\,b^{3}$,
охватывающем значения для алюминиевых сплавов [18], с центральным значением
$V^{*} = 70\,b^{3}$, измеренным для Al--Mg--Si [17]; из настоящих расчётов
$V^{*}$ не извлекался (раздел~3.3), поэтому интервал --- диапазон
чувствительности внутри модели, а не измеренная входная величина.

Результат сведён в табл.~\ref{tab:bridge}. Второй столбец получен приравниванием
левой части~(2) к 0{,}25 и решением относительно $\tau_m$:
8{,}4--62{,}7~МПа во всём интервале $V^{*}$, причём большие значения отвечают
меньшим активационным объёмам. Хотя при фиксированном $\tau_m$ усиление
пропорционально $f$, само $\tau_m$ находится обращением нелинейной функции и
потому не масштабируется как $1/f$: прямой пересчёт даёт 9{,}8--73{,}3~МПа
при $f = 0{,}001$ и 8{,}1--60{,}3~МПа при $f = 0{,}003$.

Далее следуют три сопоставления. Если бы 147~МПа работы [5] были
разрешённым касательным напряжением, действующим на дислокации,
соотношение~(2) предсказало бы усиление от $6{,}8\cdot10^{2}$ до
$2{,}7\cdot10^{46}$ (третий столбец), то есть по меньшей мере в
$2{,}7\cdot10^{3}$ раза больше наблюдаемых 0{,}25; 147~МПа в роли
касательного напряжения несовместимы с наблюдаемым усилением. Напряжение,
которое включение, деформированное на 0{,}194\%, создаёт в действительности,
меньше. Для сцепленной сферы, удерживаемой при этой деформации, точное
решение [7] даёт 41~МПа разрешённого сдвига внутри частицы и непосредственно у её поверхности
(раздел~3.1); при такой амплитуде соотношение~(2) даёт 0{,}045 при
$V^{*} = 19\,b^{3}$, 0{,}31 при $30\,b^{3}$ и значения, далеко превышающие
0{,}25, при больших $V^{*}$ (четвёртый столбец). Гребень ячейки границы,
удерживаемый при той же деформации, создаёт непосредственно над своей
вершиной 15~МПа разрешённого сдвига (раздел~3.1); при такой
амплитуде оценка даёт 0{,}042 при $V^{*} = 50\,b^{3}$ и
0{,}16 при $70\,b^{3}$ --- активационном объёме, измеренном для
Al--Mg--Si [17] (последний столбец). При деформации 0{,}194\%, таким
образом, напряжение вокруг включений имеет тот порядок, которого требует
оценка: наблюдаемые 0{,}25 воспроизводятся при $V^{*}$ примерно от
30 до 75$\,b^{3}$ в зависимости от того, какая из двух
амплитуд взята.

Это согласованность, а не подтверждение, по трём причинам. Оценка
экспоненциально чувствительна к произведению $V^{*}\tau_m$: удвоение $V^{*}$
при фиксированной амплитуде сдвигает предсказание на два--четыре порядка,
так что согласие фиксирует немногим больше порядка величины. Деформация
0{,}194\% --- входная величина, взятая из прежней оценки межфазного
напряжения [5], а не измеренное свойство включений; магнитострикция,
измеренная для объёмных сплавов железо--алюминий, не превышает $10^{-4}$
[11,\,12], то есть в двадцать раз меньше, и, поскольку напряжение
пропорционально деформации, включение с деформацией $10^{-4}$ создало бы на
своей поверхности не более 2~МПа, для которых соотношение~(2) даёт усиление
менее 0{,}4\% при любом $V^{*}$ из интервала. Наконец, оценка описывает
поле, существующее, пока включение деформировано; она ничего не говорит о
десятках минут, в течение которых эффект сохраняется после снятия поля
(раздел~5.3). Деформируются ли включения этого сплава на 0{,}194\% в поле
0{,}7~Тл --- вопрос, на котором упругий механизм стоит или падает, и вопрос
этот экспериментальный.

\begin{table}[!htbp]
\centering
\caption{Двухмасштабная оценка~(2) при $T = 300$~К, $f = 0{,}00246$ и
$b = 2{,}8638$~\AA{}. Второй столбец --- разрешённое касательное напряжение
$\tau_m$ на поверхности частицы, воспроизводящее измеренный прирост
ползучести 25\%. Остальные столбцы --- усиление
$\langle\dot\varepsilon\rangle/\dot\varepsilon_0 - 1$, которое
соотношение~(2) предсказывает, если бы в роли $\tau_m$ выступали 147~МПа
(прежняя оценка), 41~МПа (значение внутри аналитической сферы, удерживаемой при 0{,}194\%,
оно же --- амплитуда её дальнего поля, спадающего как обратный куб; 41{,}45
без округления) или 15~МПа (максимум на оси гребня в
ячейке границы, 15{,}45 без округления, раздел~3.1). Наблюдаемое значение --- 0{,}25.}
\label{tab:bridge}
\begin{tabular}{lcccc}
\hline
$V^{*}$ & $\tau_m$ для $+25\%$, МПа & при 147~МПа & при 41~МПа & при 15~МПа \\
\hline
$19\,b^{3}$  & 62{,}7 & $6{,}8\cdot10^{2}$  & 0{,}045 & 0{,}004 \\
$30\,b^{3}$  & 39{,}7 & $3{,}9\cdot10^{6}$  & 0{,}31 & 0{,}010 \\
$50\,b^{3}$  & 23{,}8 & $3{,}9\cdot10^{13}$ & 17 & 0{,}042 \\
$70\,b^{3}$  & 17{,}0 & $4{,}8\cdot10^{20}$ & $1{,}2\cdot10^{3}$ & 0{,}16 \\
$100\,b^{3}$ & 11{,}9 & $2{,}4\cdot10^{31}$ & $9{,}3\cdot10^{5}$ & 1{,}2 \\
$142\,b^{3}$ & 8{,}4  & $2{,}7\cdot10^{46}$ & $1{,}2\cdot10^{10}$ & 31 \\
\hline
\end{tabular}
\end{table}

\section{Обсуждение}

\subsection{Три напряжения и природа магнитной фазы}

В работе фигурируют три напряжения, и их необходимо различать. Первое ---
147~МПа работы [5]. Это оценка напряжения на поверхности включения,
полученная из индукции поля через эмпирический коэффициент; это не
измеренное напряжение, не разрешённое касательное напряжение в какой-либо
плоскости скольжения и не результат настоящих расчётов. Будь оно касательным
напряжением в плоскости скольжения, оно превышало бы 95--105~МПа,
при которых отрывается нижний партнёр дислокационной пары раздела~3.2;
поэтому нагружаемые ячейки не могут проверить эту оценку сравнением с
порогом --- они проверяют, создаёт ли деформированное включение в матрице
что-либо такого порядка.
Второе --- то, что даёт упругое включение, деформированное на 0{,}194\% и
сцепленное с матрицей. Точное решение для сцепленной сферы [7] --- без
проскальзывания по границе: атомы по обе её стороны взаимодействуют через
один и тот же потенциал, и никакого искусственного ограничения на границе
нет, --- даёт 41~МПа разрешённого сдвига внутри частицы со спадом
$(a/r)^{3}$ снаружи (рис.~\ref{fig:rss}). Третье --- поле, измеренное в
ячейке при той же удерживаемой деформации (раздел~3.1):
15~МПа разрешённого сдвига непосредственно над вершиной гребня
со спадом в пределах 80~\AA{} до $2{,}5\pm0{,}5$~МПа --- разрешающей
способности расчёта; в
среднем по ширине ячейки максимум составляет 5{,}0~МПа,
поскольку гребень занимает меньше половины этой ширины. Второе и третье
напряжения одного порядка и вместе ограничивают амплитуду, подставляемую в
соотношение~(2). Первое превышает их в 3{,}6--10 раз, потому что
$E\varepsilon$ --- напряжение стержня, удерживаемого при такой деформации, а
не включения, окружённого матрицей, которая деформируется вместе с ним.

Отдельная трудность связана с самой фазой. В исследованной партии сплава
рентгеновская дифракция определяет включения как Al$_{13}$Fe$_4$ (в
работе~[5] та же фаза записана как Fe$_4$Al$_{13}$; здесь всюду используется
обозначение Al$_{13}$Fe$_4$), а намагниченность образцов обнаруживает петлю
гистерезиса, то есть включения содержат ферромагнитную составляющую [5]; в
более ранних работах серии состав включений определён как Fe$_x$Al$_{1-x}$ с
$x = 86$--$87$~ат.\% [3,\,6]. Стехиометрический Al$_{13}$Fe$_4$, однако, не
обнаруживает магнитного упорядочения вплоть до 2~К --- ни по намагниченности,
ни по мёссбауэровским спектрам, непосредственно зондирующим магнитное
состояние атомов железа [19,\,20]: идеальное соединение не намагничивается,
и связанного с намагничиванием изменения формы у него быть не может (строго
нулевой деформации любого происхождения это, впрочем, не гарантирует).
Измеренный ферромагнетизм включений принадлежит, следовательно, не идеальной
решётке Al$_{13}$Fe$_4$, а иной их составляющей --- например, обогащённым
железом областям того типа, что отмечены в [3,\,6], --- и именно деформация
этой составляющей в поле существенна для механизма. Для включений данного
сплава эта деформация не измерена; использованные выше $10^{-4}$ ---
наибольшее значение, приведённое для массивных сплавов железо--алюминий
[11,\,12]. В расчётной ячейке включение представлено Al$_{13}$Fe$_4$ как
фазой, определённой в данной партии, с известной структурой [10];
деформация, заданная ему, принята как данное. Измерение деформации реальных
включений в поле представляется решающим экспериментальным шагом.

\subsection{О неосуществимости прямого моделирования и о том,
что МД измерить может}

Дислокация ведёт себя как линия, находящаяся под натяжением. Чтобы прогнуть
её через напряжённую область размером $L$, напряжение должно превышать
примерно $\mu b/L$, где $\mu$ --- модуль сдвига, $b$ --- вектор Бюргерса
[18]. Если бы все заявленные 147~МПа действовали как разрешённое касательное
напряжение --- предположение, наиболее благоприятное для механизма, --- этот
критерий даёт $L \approx 52$~нм (около $180\,b$), то есть ячейку примерно из
$2\cdot10^{8}$ атомов с учётом окружающей матрицы --- против примерно
$10^{5}$ атомов в ячейках настоящей работы. При 41~МПа соответствующий
размер составляет около 0{,}2~мкм, при 15~МПа --- около 0{,}5~мкм.
Прямое моделирование механизма на атомном масштабе, таким образом,
неосуществимо. Молекулярная динамика способна дать величины, которыми
оперирует описание на большем масштабе: напряжение на границе и его затухание
с расстоянием (раздел~3.1), напряжения, при которых дислокации приходят в
движение или образуются (разделы~3.2 и~3.3), и, через соотношение~(2),
проверку согласованности всей цепочки.

\subsection{Протокол магнитной памяти}

В экспериментах поле во время ползучести отсутствует. Образец выдерживается в
поле 30~мин, поле снимается, и лишь затем прикладывается нагрузка [4,\,6];
измеряется, таким образом, память о выдержке.
Всё, что делает поле, должно поэтому сохраняться десятки минут после его
снятия. Упругая деформация включения на это не способна: с исчезновением поля
включение возвращается к исходной форме, и напряжение вокруг него исчезает за
наносекунды. Настоящие расчёты содержат включение, удерживаемое при своей деформации,
пока релаксирует матрица (ячейка границы), включение, релаксировавшее из
него (свободный вариант), и нагружаемые ячейки со свободным (недеформированным) или удерживаемым
(недеформированным и деформированным) включением; ни одна из них не представляет состояния
материала после снятия поля. Десятки минут
способна сохраняться медленная перестройка атомов --- например, атомов Mg и
Si, закрепляющих дислокации (раздел~3.3), --- запущенная, пока действовало
поле. Такая перестройка в настоящей работе не установлена. Для её обоснования
потребовались бы термическая история и содержание вакансий в образцах, оценка
коэффициентов диффузии и отдельный расчёт, в котором поле включается,
выдерживается, выключается и лишь затем прикладывается нагрузка. Поэтому она
приводится здесь как одна из проверяемых гипотез, а не как результат.

\subsection{Ограничения}

Перенос полученных чисел на реальную частицу ограничивают два обстоятельства
геометрии. В ячейке включение --- гребень: выступ полуэллиптического сечения,
протяжённый вдоль $y$ через всю ячейку (раздел~2.1), а не компактная частица.
Напряжение вне такого выступа спадает с расстоянием медленнее, чем
$(a/r)^{3}$ для компактной частицы, а его дальность --- около 30~\AA{} в полную силу и 80~\AA{} в целом
(раздел~3.1) --- задаётся сечением гребня: это свойство данной ячейки, а не микронного включения, поле
которого простиралось бы соответственно дальше. Соотношение~(2) берёт
пространственную статистику из трёхмерного решения для компактной частицы, а
амплитуду --- либо из того же решения (41~МПа), либо из ячейки с гребнем
(15~МПа); эти значения различаются примерно в
2{,}7 раза, что и есть влияние геометрии, и это различие
проходит через табл.~\ref{tab:bridge} как разница двух её последних
столбцов. Двумерный упругий расчёт для одного гребня с теми же деформацией и
упругими постоянными, но в неограниченной матрице той же жёсткости, даёт
20~МПа у поверхности гребня со спадом до 5~МПа в пределах
10~\AA{}; атомистический гребень, жёстко удерживаемый на жёстком слое, даёт
меньшее значение у поверхности и более пологий профиль. Атомистическое и
континуальное описания гребня, таким образом, согласуются по величине, а
отличие от сферы --- геометрическое, а не артефакт ячейки.

Межатомный потенциал [9] выбран из-за охвата Al, Fe, Mg и Si; насколько точно
он воспроизводит границу Al/Al$_{13}$Fe$_4$, напряжение образования
дислокаций, строение ядра дислокации и закрепление атомами Mg и Si, в
настоящей работе отдельно не проверялось, поэтому абсолютные пороги следует
считать модельно зависимыми. Остаточное несоответствие двух решёток одинаково
в контрольной и деформированной ячейках и сокращается в их разности в первом
приближении; нелинейная релаксация может препятствовать точному сокращению.

Пороги раздела~3.2 --- наблюдения в одной ячейке, с одним набором начальных
скоростей атомов и одним темпом нагружения; их следует повторить с
независимыми наборами начальных скоростей и по меньшей мере при втором темпе.
Результат по закреплению опирается на два случайных расположения атомов Mg
и Si, лишь одно из которых прошло всю лестницу, и по одному набору начальных
скоростей; статистическая оценка порога закрепления
потребовала бы трёх--пяти независимых расположений и наборов скоростей. По
этой же причине закрепление приводится как наблюдение, а активационный объём
из него не извлекается. Термостат --- алгоритм, удерживающий температуру
ячейки при 300~К, --- действует на все подвижные атомы и сам может влиять на
пороги и на всплеск температуры при встрече и взаимном уничтожении двух
дислокаций; динамический участок следует
повторить либо с термостатом, выключенным после установления температуры,
либо с термостатом, действующим только на атомы у границ ячейки.

Поле представлено только механически --- деформацией включения. Механизмы, в
которых поле действует непосредственно на дислокацию или на препятствие, как
в теориях [1,\,2], остаются за рамками расчёта. Стоит отметить, что изменения
высоты барьера, которыми оперируют эти теории, 0{,}01--0{,}1~эВ, через
произведение $V^{*}\Delta\tau$ отвечают 0{,}97--9{,}7~МПа при
$V^{*} = 70\,b^{3}$ и 0{,}48--35{,}9~МПа во всём интервале $V^{*}$ ---
перекрываясь с найденными здесь напряжениями, а не совпадая с ними. Сами
напряжения вычислены из сил и положений атомов (вириальное выражение [8]) с
усреднением по слоям толщиной 4~\AA{}, параллельным границе; это приближение,
разрешающая способность которого после усреднения шести компонент тензора
по слою составляет $2{,}5 \pm 0{,}5$~МПа. Наконец, соотношение~(2) --- отдельная
модель, связывающая локальное напряжение с изменением скорости ползучести;
зависимости самих МД-порогов от темпа нагружения оно не устраняет.

\section{Заключение}

Напряжение, создаваемое деформированным включением в матрице, рассчитано в
ячейке границы Al/Al$_{13}$Fe$_4$ из 91\,428 атомов при удлинении включения
вдоль направления поля на 0{,}194\% без изменения объёма --- деформации,
отвечающей прежней оценке 147~МПа, --- при удерживаемой деформации
включения. Разрешённое касательное напряжение в матрице достигает
15~МПа непосредственно над включением и спадает до уровня
дальнего поля $2{,}5\pm0{,}5$~МПа в пределах 80~\AA{}; в среднем по ширине
ячейки максимум составляет 5{,}0~МПа. Аналитическое решение
для компактной частицы, удерживаемой при той же деформации, даёт 41~МПа
внутри неё со спадом обратно пропорционально кубу расстояния, так что для
включения микронного размера то же поле простирается на доли микрона.
Включение, которое ничто не удерживает, снимает большую часть деформации.

В нагружаемых ячейках заранее введённая пара дислокаций рядом с гребнем
разрывается при приложенном сдвиге 95--105~МПа при свободно релаксирующем включении; при удерживаемом включении деформированная и недеформированная ячейки на каждом этапе совпадают в пределах одного кадра рампы, 9~МПа; новых дислокаций на границе в недеформированной ячейке не образуется вплоть до конца рампы при 400~МПа; а случайное
расположение атомов Mg и Si удерживает дислокацию вплоть до 75~МПа.
Первые пикосекунды каждой рампы, при нулевом приложенном напряжении, отвечают на вопрос, сдвигает ли одно лишь поле деформированного гребня помещённую рядом дислокацию: один и тот же партнёр движется одинаково при деформированном и недеформированном гребне, движимый полем когерентности границы, а не деформацией гребня. Эти значения относятся к выбранным ячейкам, темпу нагружения и
потенциалу и не являются константами материала.

Двухмасштабная оценка~(2) требует разрешённого касательного напряжения
8{,}4--62{,}7~МПа на поверхности включения для повышения скорости
ползучести на 25\% при $f = 0{,}00246$ и $V^{*} = 19$--$142\,b^{3}$.
Рассчитанные амплитуды --- 41~МПа для частицы и 15~МПа для
гребня --- воспроизводят наблюдаемое усиление при $V^{*}$ от 30 до
75$\,b^{3}$, то есть в интервале, содержащем значение, измеренное для
Al--Mg--Si; 147~МПа в роли разрешённого сдвига превысили бы наблюдаемое
усиление не менее чем в $2{,}7\cdot10^{3}$ раза. Упругий механизм,
следовательно, количественно состоятелен тогда и только тогда, когда
включения деформируются в поле примерно на 0{,}2\%: магнитострикция,
измеренная для объёмных сплавов железо--алюминий, в двадцать раз меньше и
дала бы усиление менее 0{,}4\%.

Поскольку во время испытания на ползучесть поле выключено, его действие
должно сохраняться десятки минут; упругая деформация на это не способна, и в
качестве оставшейся возможности указывается медленная перестройка атомов ---
гипотеза, настоящими расчётами не проверенная. Решающими дальнейшими шагами
представляются установление ферромагнитной составляющей включений, прямое
измерение их деформации в поле и расчёт, в котором поле включается,
выдерживается, выключается и лишь затем прикладывается нагрузка.
